(1) 随物质导数 \(D_t=\partial_t+U\partial_x\)。质量元的横向加速度是 \(D_t^2y\)，张力横向合力每长度为 \(Ty_{xx}\)，故 \(\boxed{\mu(\partial_t+U\partial_x)^2y=Ty_{xx}}\)，定义相对带的横波速度 \(c=\sqrt{T/\mu}\)。

(2) 取 \(y=Ae^{i(kx-\omega t)}\) 得 \((\omega-Uk)^2=c^2k^2\)，故 \(\boxed{\omega_\pm=(U\pm c)k}\)、\(v_{g,\pm}=U\pm c\)。当 \(|U|<c\) 时一个波包逆着实验室正方向传播；当 \(U>c\) 时两支在实验室均向下游移动，信息不能逆流到上游。

(3) 令物质横向速度 \(v=D_ty\)，方程写 \(\mu D_tv=Ty_{xx}\)。乘以 \(v\) 得 \(D_t[\mu v^2/2+Ty_x^2/2]+\partial_x(-Tvy_x)=0\)，即对实验室变量有 \(\partial_t e+\partial_x(Ue-Tvy_x)=0\)。前一项为随带能量对流，后一项为相对带的弹性波能流。

(4) 令 \(U\to0\) 就得到 \(\mu y_{tt}=Ty_{xx}\) 和波速 \(\pm c\) 的静止弦。\(U\) 仅作Galilean平移，对各本征波产生 \(Uk\) 的多普勒位移；若驱动带的电机速度恒定，局部波动自身无负阻尼，不可仅因下游群速变大就称出现激光式增益。
